\chapter{问题定义与设计目标}

\section{问题定义}

给定目标 ISA、执行后端、待扫描编码空间和当前运行环境，本文的目标是判定每条编码在“处理器执行结果”与“反汇编器合法性判断”之间是否存在行为分歧，并输出可复现、可审计的日志与元数据。这里的“分歧”不只包括狭义的 hidden instruction，也包括特权相关拒绝、执行期异常、超时和工具链之间的不一致。

本文采用的工作假设是：在用户态可控环境中，单条指令注入、信号捕获和反汇编判断能够构成一个实用的差异审计闭环。这个闭环并不试图直接给出 authoritative semantics，而是提供一组可复核的异常候选，并把后续 triage 的上下文一并保留下来。

\section{分类口径}

当前系统使用如下工作分类。
\begin{itemize}[leftmargin=2em]
  \item \textbf{H}：执行成功，但主反汇编器不认为其为合法指令。
  \item \textbf{D}：主反汇编器认为合法，但执行在测试指令处被处理器拒绝。
  \item \textbf{T}：执行超时。
  \item \textbf{X}：处理器进入执行路径后，在执行期间触发其他异常。
  \item \textbf{P}：特权或权限相关行为。
  \item \textbf{UNKNOWN}：当前实验无法稳定归类，或仍需人工 triage。
\end{itemize}

这些标签在论文叙事中已经足够表达问题空间，但在当前实现中，它们并不是完全等价的自动输出标签。尤其是 AArch64 的 controlled illegal 与 trap 路径，已经通过受控实验暴露出日志 taxonomy 的空洞。因此，本文在正文中显式区分“理论上希望表达的分类口径”和“当前程序已经稳定输出的 artifact 标签”，避免把实现边界隐藏到附录中。

\begin{table}[t]
  \centering
  \caption{当前工作分类与日志自动输出能力}
  \label{tab:taxonomy}
  \begin{tabular}{p{1.5cm}p{4.0cm}p{2.5cm}p{5.0cm}}
    \toprule
    标签 & 含义 & 当前自动输出 & 当前已知缺口 \\
    \midrule
    H & 执行成功，但主反汇编器不认为合法 & 部分是 & RISC-V 已有 raw/filter 双口径；AArch64 当前样例尚未覆盖稳定 H \\
    D & 主反汇编器认为合法，但执行在测试指令处被拒绝 & 是 & AArch64 privileged-in-user 当前折叠到 D(strict) \\
    T & 执行超时 & 是 & 仅对应当前 watchdog/timeout 口径 \\
    X & 执行期间触发其他异常 & 部分是 & 已知 known-disassembly exec-fault 路径未稳定进入 artifact \\
    P & 特权或权限相关行为 & 否 & 当前更多依赖人工 triage，而非独立自动标签 \\
    UNKNOWN & 当前无法稳定归类 & 否 & 需结合第二解码器与人工复核产生 \\
    \bottomrule
  \end{tabular}
\end{table}

\section{覆盖模型与边界}

为了避免过度主张，本文在问题定义阶段就冻结三个边界。第一，扫描结果总是绑定到具体 ISA、具体执行后端、具体编码子空间和具体运行环境；脱离这些条件的泛化结论都不成立。第二，Layer A 的模拟或容器结果不能直接替代真实硬件结论。第三，用户态审计天然看不到全部微架构状态，因此它适合发现差异候选和实现偏差，但不等价于完整的 ISA 合规证明。

\section{设计目标}

本文的系统设计围绕以下目标展开。
\begin{enumerate}[leftmargin=2em]
  \item 单条可控执行：每次只测试一个目标编码，并能在测试后回收控制流。
  \item 程序状态可回收：局部异常不应破坏整个 harness。
  \item 用户态可部署：尽量避免依赖内核修改或特权模块。
  \item 多 ISA 共用编排骨架：前端逻辑不因架构差异而大规模分叉。
  \item 新 ISA 接入成本可量化：明确区分可自动生成部分与必须人工加固部分。
  \item 容器或仿真结果能够与后续实机结果对照。
\end{enumerate}

这六个目标分别对应本文后续章节中的设计选择：第 4 章解释统一架构与双后端如何实现 G1--G4，第 5 章解释生成器如何回答 G5，第 6 章则用 Layer A/B/C 结构回应 G6。
